% Figure: the modelling cycle. Six labelled stages around a loop, with a
% feedback arrow from validation back to abstraction (the iterative
% revision the chapter installs as discipline).
\begin{figure}[htbp]
\centering
\begin{tikzpicture}[
  stage/.style={draw=engblack, fill=engblue!12, rounded corners=3pt,
    minimum width=2.6cm, minimum height=0.85cm, font=\small\bfseries,
    align=center},
  arrow/.style={-{Stealth}, thick, draw=engblack!70},
  feedback/.style={-{Stealth}, thick, dashed, draw=engred!80},
  node distance=0.4cm and 0.6cm,
]
% Six stages around a hexagon
\node[stage] (q)  at ( 0.0,  3.0) {Question};
\node[stage] (a)  at ( 4.0,  1.8) {Abstraction};
\node[stage] (i)  at ( 4.0, -1.8) {Identification};
\node[stage] (v)  at ( 0.0, -3.0) {Validation};
\node[stage] (u)  at (-4.0, -1.8) {Uncertainty};
\node[stage] (l)  at (-4.0,  1.8) {Look-back};

% Forward cycle
\draw[arrow] (q) -- (a);
\draw[arrow] (a) -- (i);
\draw[arrow] (i) -- (v);
\draw[arrow] (v) -- (u);
\draw[arrow] (u) -- (l);
\draw[arrow] (l) -- (q);

% Feedback: validation forces revised abstraction
\draw[feedback] (v) to[bend left=25] node[midway, fill=white, inner sep=2pt,
  font=\scriptsize\itshape, align=center] {revise on\\ residual\\ structure} (a);

% Centre note
\node[font=\small\itshape, align=center, color=engblack!70] at (0, 0)
  {the modelling cycle\\ (one iteration)};
\end{tikzpicture}
\caption[The modelling cycle]{The modelling cycle as the chapter
installs it. The solid arrows mark one full pass: question, abstraction,
identification, validation, uncertainty report, look-back. The dashed
arrow marks the iterative revision triggered when validation reveals
residual structure that the current abstraction does not capture; the
revised model re-enters the cycle at the abstraction stage. A model
that has not closed the loop at least once has not yet been validated;
a model that has closed the loop without revision has either been
correct on the first abstraction (rare) or has been built without
inspecting the residuals (common, and the failure mode of
section~18.8).}
\label{fig:vol02:ch18:modelling-cycle}
\end{figure}
